\section{MCP + skills registry}

\begin{frame}[fragile]{MCP: making the stack portable}
  \cmd{eda-mcp} (stdio Model Context Protocol server) exposes the repo's
  primitives to any MCP-aware client (Claude Code, IDE extensions, custom
  tools):
  \begin{itemize}\footnotesize
    \item \term{Skills}: named prompts rendered from Markdown bundles
          (\cmd{analog.gmid\_sizing}, \cmd{digital.cocotb\_testbench},
          \cmd{flow.drc\_checker}, \ldots).
    \item \term{LUTs}: gm/ID tables for all registered PDK/device pairs.
    \item \term{Runners}: SPICE, LibreLane, DRC, LVS -- same code the
          agents call internally.
    \item \term{Topology registry}: list + metadata of circuit topologies.
  \end{itemize}
  \begin{lstlisting}[style=shellstyle]
# Register the server in ~/.config/claude-code/mcp.json
{"mcpServers": {"eda-agents": {"type": "stdio",
                "command": "eda-mcp"}}}
  \end{lstlisting}
  \footnotesize
  Full example: \cmd{src/eda\_agents/templates/mcp.json}.
  \note{%
    The MCP server is the integration seam. Without it, every consumer
    has to depend on the \cmd{eda\_agents} Python package and know the
    internal module layout. With it, the consumer asks for
    \cmd{get\_skill(\"analog.gmid\_sizing\")} or
    \cmd{run\_librelane(project\_dir, to=\"SIGNOFF\")} by name.\par
    Namespaces shipped today: \cmd{analog.*} (10+ skills),
    \cmd{digital.*} (cocotb TB authoring, RTL-to-GDS driver, etc.),
    \cmd{flow.*} (DRC checker, DRC fixer). The \cmd{list\_skills}
    tool enumerates them at runtime so an agent can discover what is
    available.\par
    LUT access is non-trivial: IHP SG13G2 LUTs live in an external
    \cmd{ihp-gmid-kit} clone; GF180 LUTs auto-download on first use and
    cache under \cmd{\char126/.cache/eda-agents/gmid\_luts/}. The MCP
    server hides all of that.}
\end{frame}
